\section{Live Runtime Experiments}
\label{sec:livexp}
This section tests the runtime of Section~\ref{sec:live}: that delay and memory stay bounded over a long run, and that a failing camera is noticed quickly and without false alarms. Recorded KAIST test sequences are replayed at real speed as the two cameras. The runtime uses CFFM-Net v2 trained on KAIST with the memory of Section~\ref{sec:kaistmemory} and ByteTrack. It ran on the same \latLocalDevice{} as Supplementary Table~S11; its inference times are higher because they include the memory, pre-processing and the tracker.

\subsection{Soak test}
\label{sec:soak}
Paired \cfgLatW{}$\times$\cfgLatH{} cameras ran at 30 frames per second for ten minutes (\livSoakThirtySeconds{}\,s) and at 20 frames per second for \cfgSoakShortSeconds{}\,s.

Supplementary Table~S9 gives the full soak and link-fault measurements. The stability criterion is met: over the ten-minute run the end-to-end delay drifts by \livSoakThirtyDriftMs{}\,ms and GPU memory changes by \livSoakThirtyGpuGrowthMb{}\,MB, so neither grows. The runtime is not real time at the camera rate on this hardware, however. At a nominal 30 frames per second it processes \livSoakThirtyFps{} and drops \livSoakThirtyDropPercent{}\,\% of the captured frames, and at 20 frames per second it processes \livSoakTwentyFps{} and drops \livSoakTwentyDropPercent{}\,\%. The median inference time with memory and tracking is \livSoakThirtyInferMedianMs{}\,ms and the median end-to-end delay \livSoakThirtyEndToEndMedianMs{}\,ms at 30 frames per second, with a 95th percentile of \livSoakThirtyEndToEndPNinetyFiveMs{}\,ms. The 20 frames per second run drifts by \livSoakTwentyDriftMs{}\,ms, about eight percent of its median delay, and is too short to show whether that settles.

\subsection{Link faults and the stream monitor}
\label{sec:linkfaults}
A scripted run of \cfgLinkSeconds{}\,s at 20 frames per second applies three faults: the thermal cable pulled for \cfgUnplugSeconds{}\,s from second 20, the visible camera frozen for \cfgUnplugSeconds{}\,s from second 50, and the thermal stream stalled for \cfgStallSeconds{}\,s from second 80. Throughout, \cfgLinkDropPercent{}\,\% of frames are dropped at random and the thermal stream is delayed by \cfgLinkDelayMs{}\,ms. We repeat the run with the monitor feeding the sensor flags to the model (on) and with the monitor observing only (off).

Supplementary Table~S8 gives the monitor timing for each fault. A cable pull is flagged \livFaultOnUnpluggedDetectS{}\,s after onset and cleared \livFaultOnUnpluggedRecoverS{}\,s after the camera returns, a stall is flagged after \livFaultOnStalledDetectS{}\,s and cleared after \livFaultOnStalledRecoverS{}\,s, and a frozen visible camera is flagged after \livFaultOnFrozenDetectS{}\,s, which is longer because a freeze is recognised only after consecutive nearly identical frames and a replayed stream drops frames at random. The monitor is not free of false alarms: the thermal stream is marked down outside the scripted windows in \livFaultOnUnpluggedFalseAlarmPercent{}\,\% of records in the cable-pull measurement and \livFaultOnStalledFalseAlarmPercent{}\,\% in the stall measurement, which we attribute to the random drops and the added delay but did not test separately. The visible camera had none.

Feeding the flags to the model gives no measured accuracy benefit. The miss rate over all frames is \livFaultOnMr{}\,\% with the monitor on and \livFaultOffMr{}\,\% with it off, a difference of \livFaultMrChange{}\,points against the monitor, and throughput and delay are the same within a fraction of a frame (\livFaultOnFps{} against \livFaultOffFps{} frames per second). This is one run, the scripted faults cover only fifteen of the \cfgLinkSeconds{} seconds, and the memory-trained detector may already cope with a black thermal frame, in line with the flagged-versus-unflagged result of Section~\ref{sec:faults}. We conclude that the monitor detects and clears link faults quickly, and that we have not shown that its flags improve detection.

\paragraph{Supplementary material}
Supplementary video S2 (\texttt{live\_session.mp4}, about nine seconds at 20 frames per second) shows the live board during a thermal cable pull.
